% Shared notation (paper/BRIEF.md §3). Use these macros; do not redefine them in a section.
\makeatletter
% \X_{..} after a macro that already carries a subscript joins the two: \Mdiv_{ij} -> M_{div,ij}.
\newcommand{\@joinsub}[2]{\@ifnextchar_{\@joinsub@{#1}{#2}}{#1_{#2}}}
\def\@joinsub@#1#2_#3{#1_{#2,#3}}
\makeatother
\newcommand{\R}{\mathbb{R}}
\newcommand{\Omegac}{\Omega}                       % cavity volume
\newcommand{\bdry}{\partial\Omega}                 % PEC wall
\newcommand{\E}{\mathbf{E}}
\newcommand{\Hf}{\mathbf{H}}
\newcommand{\curl}{\nabla\times}
\newcommand{\Hcurl}{H_0(\mathrm{curl};\Omega)}
\newcommand{\Vh}{V_h}                              % discrete H(curl) space (Nedelec, order p, curved mesh)
\newcommand{\Sh}{S_h}                              % discrete H^1_0 potential space, order p+1
\newcommand{\Kmat}{\mathsf{K}}                     % curl-curl stiffness
\newcommand{\Mmat}{\mathsf{M}}                     % mass
\newcommand{\Gmat}{\mathsf{G}}                     % discrete gradient (potentials -> H(curl) DOFs)
\makeatletter\DeclareRobustCommand{\Kp}{\@joinsub{\mathsf{K}}{p}}\makeatother                     % G^T M G
\newcommand{\Pmat}{\mathsf{P}}                     % kernel projection
\newcommand{\Vmat}{\mathsf{V}}                     % learned basis in DOFs, n x m
\newcommand{\Cmat}{\mathsf{C}}                     % fixed point-values -> DOF map
\makeatletter\DeclareRobustCommand{\Mdiv}{\@joinsub{\mathsf{M}}{\mathrm{div}}}\makeatother
\makeatletter\DeclareRobustCommand{\Av}{\@joinsub{\mathsf{A}}{V}}\makeatother
\newcommand{\ndof}{n}
\newcommand{\nbasis}{m}                            % width of the learned basis
\newcommand{\nmodes}{K}                            % number of predicted modes
\makeatletter\DeclareRobustCommand{\lamh}{\@joinsub{\lambda}{h}}\makeatother  % discrete eigenvalue; \lamh_k -> lambda_{h,k}
\makeatletter\DeclareRobustCommand{\Qz}{\@joinsub{Q}{0}}\makeatother
\newcommand{\RoQ}{R/Q}
\makeatletter\DeclareRobustCommand{\Epk}{\@joinsub{E}{\mathrm{pk}}}\makeatother
\makeatletter\DeclareRobustCommand{\Bpk}{\@joinsub{B}{\mathrm{pk}}}\makeatother
\makeatletter\DeclareRobustCommand{\Eacc}{\@joinsub{E}{\mathrm{acc}}}\makeatother
\makeatletter\DeclareRobustCommand{\Lacc}{\@joinsub{L}{\mathrm{acc}}}\makeatother
\newcommand{\LF}{\mathrm{LF}}
\newcommand{\HF}{\mathrm{HF}}
\newcommand{\relLtwo}{\mathrm{rel\text{-}L^2}}
\DeclareMathOperator{\spn}{span}
\DeclareMathOperator{\diag}{diag}
